\subsection{}
Deux cordes de masses linéiques différentes sont attachées ensemble et mises sous tension. Une onde sinusoïdale progressive est produite sur la corde la plus légère. Que se passe-t-il quand cette onde est transmise à la deuxième corde (plus lourde) ? (On s'intéresse juste à la partie transmise) Justifiez.

\begin{enumerate}[label=\Alph*)]
    \item La fréquence et la longueur d'onde augmentent
    \item La fréquence est inchangée et la longueur d'onde augmente
    \item La fréquence est inchangée et la longueur d'onde diminue
    \item La fréquence et la longueur d'onde diminuent
    \item La fréquence augmente et la longueur d'onde est inchangée
    \item La fréquence diminue et la longueur d'onde est inchangée
\end{enumerate}

\noindent\textbf{Justification :} \hrulefill \\[2.5ex]
\rule{\linewidth}{0.4pt} \\[2.5ex]
\rule{\linewidth}{0.4pt}
\vfill